%%%%%%%%%%%%%%%%%%%%%%%%%%%%%%%%%%%%%%%%%%%%%%%%%%%%%%%%%%%%%%%%%%%%%
%% sec_theory_EM_v2.tex
%%
%% Sections 2.2 -- 2.4 for paperH_HLmain.tex
%% RMG Periodic RT-TDDFT -- Paper 2
%% REVISED April 2026: clean separation of EM-field perturbations
%% (gauge-dependent, Track A) from scalar-potential perturbations
%% (gauge-independent, Track B).
%%
%% Subsection structure:
%%   2.2  Interaction with electromagnetic radiation
%%        2.2.1  Plane wave and minimal coupling
%%        2.2.2  Multipole expansion and long-wavelength approximation
%%        2.2.3  Crystal momentum conservation: vertical transitions
%%        2.2.4  Beyond the optical limit: scalar-potential perturbations
%%   2.3  Gauge transformation: length to velocity gauge
%%        2.3.1  Why the length gauge fails for periodic systems
%%        2.3.2  The Goppert-Mayer transformation
%%        2.3.3  The velocity-gauge TDKS Hamiltonian
%%        2.3.4  The perturbative momentum kick
%%   2.4  Vector potential and magnetic perturbations
%%        2.4.1  Electric perturbation: spatially uniform A(t)
%%        2.4.2  Magnetic perturbation: spatially varying A(r)
%%        2.4.3  Scope of the present implementation
%%
%% CONVENTIONS (all from notation.tex):
%%   Atomic units: hbar = m_e = e = 4*pi*eps_0 = 1
%%   Electron charge = -1  (e > 0 is elementary charge magnitude)
%%   Imaginary unit: \imath  (dotless i, consistent with paper 1)
%%   \myblue{} for all annotations and placeholders -- NEVER \myred{}
%%
%% SOURCES:
%%   NEAT-notes-RMG_main.tex lines 116-716 (primary)
%%   NEAT-project-notes_main.tex lines 256-716 (supplementary)
%%   Physics discussion notes, April 2026
%%%%%%%%%%%%%%%%%%%%%%%%%%%%%%%%%%%%%%%%%%%%%%%%%%%%%%%%%%%%%%%%%%%%%


%--------------------------------------------------------------------
\subsection{Interaction of Matter with Electromagnetic Radiation}
\label{sec:em-interaction}
%--------------------------------------------------------------------

The starting point for any RT-TDDFT calculation is the coupling of
the electronic system to an external time-dependent perturbation.
Two physically distinct classes of perturbation are relevant to this
work and to its natural extensions:
\emph{electromagnetic field perturbations}, in which the coupling
enters through the vector potential in the kinetic energy operator
and whose treatment requires a choice of gauge
(Sections~\ref{sec:plane-wave}--\ref{sec:vertical}); and
\emph{scalar potential perturbations}, in which a time-dependent
external potential is added directly to the Kohn-Sham effective
potential without any gauge considerations
(Section~\ref{sec:scalar-pert}).
The present work employs the electromagnetic coupling exclusively;
the scalar perturbation framework is described briefly to place the
implementation in context and to indicate natural extensions.


%....................................................................
\subsubsection{The electromagnetic plane wave and minimal coupling}
\label{sec:plane-wave}
%....................................................................

A monochromatic electromagnetic plane wave propagating in direction
$\hat{\mathbf{q}}$ with polarization $\hat{\varepsilon}$ is described
in the Coulomb gauge ($\nabla\cdot\Avec=0$, $\phi=0$) by the
vector potential
\begin{equation}
  \Avec(\mathbf{r},t)
  = A_0\,\hat{\varepsilon}\,
    e^{\imath(\mathbf{q}\cdot\mathbf{r}-\omega t)}
    + \text{c.c.},
  \label{eq:plane-wave-A}
\end{equation}
where $|\mathbf{q}|=\omega/c$ and $A_0$ is the field amplitude.
The Coulomb gauge condition $\nabla\cdot\Avec=0$ requires
$\hat{\varepsilon}\cdot\mathbf{q}=0$: the polarization is
perpendicular to the propagation direction (transversality).
The electric and magnetic fields follow from the potentials,
\begin{align}
  \Evec(\mathbf{r},t)
  &= -\frac{\partial\Avec}{\partial t}
   = \imath\omega A_0\,\hat{\varepsilon}\,
     e^{\imath(\mathbf{q}\cdot\mathbf{r}-\omega t)}
     + \text{c.c.},
  \label{eq:plane-wave-E}
  \\
  \Bvec(\mathbf{r},t)
  &= \nabla\times\Avec
   = \imath\mathbf{q}\times(A_0\hat{\varepsilon})\,
     e^{\imath(\mathbf{q}\cdot\mathbf{r}-\omega t)}
     + \text{c.c.},
  \label{eq:plane-wave-B}
\end{align}
where in eq.~(\ref{eq:plane-wave-B}) we used
$\nabla e^{\imath\mathbf{q}\cdot\mathbf{r}}
=\imath\mathbf{q}\,e^{\imath\mathbf{q}\cdot\mathbf{r}}$.
Both $\Evec$ and $\Bvec$ carry the same spatial phase factor
$e^{\imath\mathbf{q}\cdot\mathbf{r}}$ as $\Avec$.

The coupling of the plane wave to a Kohn-Sham electron is introduced
through minimal substitution
$\hat{\mathbf{p}}\to\hat{\mathbf{p}}+\Avec(\mathbf{r},t)$,
giving the time-dependent Kohn-Sham Hamiltonian
\begin{equation}
  \hat{H}(t)
  = \frac{1}{2}
    \!\left(\hat{\mathbf{p}}+\Avec(\mathbf{r},t)\right)^{\!2}
    + V_\mathrm{KS}(\mathbf{r},t).
  \label{eq:H-minimal}
\end{equation}
Expanding the kinetic term in powers of $A_0$:
\begin{equation}
  \hat{H}(t)
  = \underbrace{
      -\frac{\nabla^2}{2}+V_\mathrm{KS}
    H_HL{\hat{HH_HL0}
  + \underbrace{
      \Avec(\mathbf{r},t)\cdot\hat{\mathbf{p}}
    H_HL{\hat{H}^{(1)},\;\mathcal{O}(A_0)}
  + \underbrace{
      \tfrac{1}{2}\Avec^2
    H_HL{\hat{H}^{(2)},\;\mathcal{O}(A_0^2)}.
  \label{eq:H-expand}
\end{equation}
In the linear response regime the diamagnetic $\Avec^2$ term is
second order in the field amplitude and is
neglected.\cite{yabana-bertsch-1996,Mattiat-Luber-2022}
The electromagnetic interaction Hamiltonian is therefore
\begin{equation}
  \hat{HH_HL\mathrm{int}(\mathbf{r},t)
  = \Avec(\mathbf{r},t)\cdot\hat{\mathbf{p}}
  = A_0\,\hat{\varepsilon}\cdot\hat{\mathbf{p}}\;
    e^{\imath(\mathbf{q}\cdot\mathbf{r}-\omega t)}
    + \text{c.c.}
  \label{eq:Hint-exact}
\end{equation}
The spatial dependence of $\hat{HH_HL\mathrm{int}}$ enters entirely
through the phase factor $e^{\imath\mathbf{q}\cdot\mathbf{r}}$,
which couples the field to the electronic system.
Because this coupling appears in the \emph{kinetic energy} operator
through minimal substitution, the treatment of the
interaction~(\ref{eq:Hint-exact}) requires a choice of gauge ---
in contrast to the scalar potential perturbations discussed in
Section~\ref{sec:scalar-pert}, which carry no gauge ambiguity.


%....................................................................
\subsubsection{The multipole expansion and the long-wavelength
  approximation}
\label{sec:multipole}
%....................................................................

Expanding the spatial phase factor in eq.~(\ref{eq:Hint-exact})
as a Taylor series about $\mathbf{q}\cdot\mathbf{r}=0$:
\begin{equation}
  e^{\imath\mathbf{q}\cdot\mathbf{r}}
  = 1
  + \imath(\mathbf{q}\cdot\hat{\mathbf{r}})
  + \tfrac{1}{2}[\imath(\mathbf{q}\cdot\hat{\mathbf{r}})]^2
  + \cdots,
  \label{eq:multipole}
\end{equation}
the matrix element of $\hat{HH_HL\mathrm{int}$ between Kohn-Sham
orbitals $|\psi_a\rangle$ and $|\psi_b\rangle$ becomes:
\begin{equation}
  \langle\psi_b|\hat{HH_HL\mathrm{int}|\psi_a\rangle
  \propto
  \underbrace{
    \langle\psi_b|\hat{\varepsilon}\cdot\hat{\mathbf{p}}|\psi_a\rangle
  H_HL{\text{E1: electric dipole, velocity form}}
  + \imath\!
  \underbrace{
    \langle\psi_b|(\mathbf{q}\cdot\hat{\mathbf{r}})
    (\hat{\varepsilon}\cdot\hat{\mathbf{p}})|\psi_a\rangle
  H_HL{\text{M1}+\text{E2: magnetic dipole + electric quadrupole}}
  + \cdots
  \label{eq:multipole-me}
\end{equation}
Each successive term carries an additional factor of $|\mathbf{q}|\,a_0$
relative to the preceding one.
For near-ultraviolet radiation at $\lambda=400$~nm,
$|\mathbf{q}|=2\pi/\lambda\approx 8\times10^{-4}\,a_0^{-1}$, so
that $|\mathbf{q}|\,a\approx 1.5\times10^{-2}$ for a unit cell of
linear dimension $a=10$~\AA.
The multipole series converges rapidly and the
\emph{long-wavelength (electric dipole) approximation}
\begin{equation}
  e^{\imath\mathbf{q}\cdot\mathbf{r}} \approx 1,
  \qquad |\mathbf{q}|\,a \ll 1,
  \label{eq:LWA}
\end{equation}
is quantitatively excellent for all optical and UV spectroscopy.
Under this approximation $\Avec(\mathbf{r},t)\to\Avec(t)$ becomes
spatially uniform, and the interaction Hamiltonian reduces to
\begin{equation}
  \hat{HH_HL\mathrm{int}(t)
  = \Avec(t)\cdot\hat{\mathbf{p}},
  \label{eq:Hint-LWA}
\end{equation}
which no longer contains $\hat{\mathbf{r}}$ and is the starting
point for the velocity-gauge formulation of Section~\ref{sec:gauge}.


%....................................................................
\subsubsection{Crystal momentum conservation: vertical transitions}
\label{sec:vertical}
%....................................................................

For a periodic crystal with Bloch states
$|\psi_{n\mathbf{k}}\rangle$, the matrix element of
$\hat{HH_HL\mathrm{int}$ contains the overlap integral of the Bloch
factors $e^{\imath\mathbf{k}\cdot\mathbf{r}}$,
$e^{\imath\mathbf{k}'\cdot\mathbf{r}}$, and
$e^{\imath\mathbf{q}\cdot\mathbf{r}}$ over the unit cell.
This integral vanishes unless $\mathbf{k}'=\mathbf{k}+\mathbf{q}$,
expressing crystal momentum conservation:
\begin{equation}
  \langle\psi_{m\mathbf{k}'}|
  e^{\imath\mathbf{q}\cdot\mathbf{r}}\hat{\varepsilon}\cdot\hat{\mathbf{p}}
  |\psi_{n\mathbf{k}}\rangle
  \propto \delta_{\mathbf{k}',\mathbf{k}+\mathbf{q}}.
  \label{eq:k-conservation}
\end{equation}
The photon deposits its wavevector $\mathbf{q}$ into the electronic
system, inducing transitions with $\Delta\mathbf{k}=\mathbf{q}$.
In the long-wavelength limit $\mathbf{q}\to\mathbf{0}$,
eq.~(\ref{eq:k-conservation}) reduces to $\delta_{\mathbf{k}',\mathbf{k}}$:
all transitions are \emph{vertical} ($\Delta\mathbf{k}=0$) in the
Brillouin zone.
For optical photons the shift $|\mathbf{q}|/(\pi/a)
\approx 10^{-3}$ is negligible on the scale of the Brillouin zone,
confirming that the long-wavelength approximation is well-controlled
and quantifiable, not an assumption of unknown accuracy.


%....................................................................
\subsubsection{Beyond the optical limit: scalar-potential perturbations
  and finite-wavevector response}
\label{sec:scalar-pert}
%....................................................................

A physically distinct class of perturbations is obtained by adding
a time-dependent potential directly to the Kohn-Sham effective
potential, without any reference to an electromagnetic field:
\begin{equation}
  V_\mathrm{KS}(\mathbf{r},t)
  \to V_\mathrm{KS}(\mathbf{r},t)
     + \delta V(\mathbf{r},t).
  \label{eq:scalar-pert}
\end{equation}
Because $\delta V$ enters as a scalar potential, no vector potential
is introduced, no gauge transformation is required, and the
discussion of Sections~\ref{sec:gauge}--\ref{sec:vector-potential}
does not apply to perturbations of this type.
The Coulomb gauge is maintained throughout; the separation of the
EM field discussion from the scalar potential discussion is exact,
not approximate.

The specific choice
\begin{equation}
  \delta V(\mathbf{r},t)
  = V_0\,e^{\imath\mathbf{q}\cdot\mathbf{r}}\,\delta(t),
  \label{eq:density-wave-kick}
\end{equation}
with $\mathbf{q}$ a free parameter independent of $\omega$,
excites a density wave at wavevector $\mathbf{q}$ and all
frequencies simultaneously.
The Fourier component of the density response,
$\delta\rho(\mathbf{q},t) =
\int_\Omega\delta\rho(\mathbf{r},t)\,e^{-\imath\mathbf{q}\cdot\mathbf{r}}
\,d^3r$,
yields the density--density response function
\begin{equation}
  \chi(\mathbf{q},\omega)
  = \frac{\widetilde{\delta\rho}(\mathbf{q},\omega)}{V_0},
  \qquad
  S(\mathbf{q},\omega) \propto \mathrm{Im}\,\chi(\mathbf{q},\omega),
  \label{eq:chi-q}
\end{equation}
after Fourier transformation in time.
Repeating for a sequence of wavevectors $\mathbf{q}$ commensurate
with the simulation cell gives access to the full dynamic structure
factor $S(\mathbf{q},\omega)$, from which plasmon dispersion curves,
interband transitions at finite momentum transfer, and quantities
measurable by inelastic X-ray scattering or electron energy loss
spectroscopy (EELS) can be extracted.\myblue{[cite: EELS and IXS
RT-TDDFT references]}

A closely related scalar perturbation is the Coulomb potential of a
point charge placed at position $\mathbf{R}$ in the simulation cell,
\begin{equation}
  \delta V_\mathrm{Coul}(\mathbf{r})
  = \frac{V_0}{|\mathbf{R}-\mathbf{r}|},
  \label{eq:Coulomb-pert}
\end{equation}
which was implemented and validated in ref.~\cite{jakowski-tddft-coulomb-2019}
\myblue{[confirm cite key in references.bib]} as a model for a focused
electron beam at position $\mathbf{R}$.
Via the Fourier expansion
$1/|\mathbf{R}-\mathbf{r}|
=(4\pi/\Omega)\sum_{\mathbf{q}\neq\mathbf{0}}
e^{\imath\mathbf{q}\cdot(\mathbf{r}-\mathbf{R})}/|\mathbf{q}|^2$,
the Coulomb perturbation is a superposition of all plane-wave
components~(\ref{eq:density-wave-kick}) weighted by $1/|\mathbf{q}|^2$:
it samples $\chi(\mathbf{q},\omega)$ at all wavevectors simultaneously
but emphasizes the long-wavelength ($|\mathbf{q}|\to0$) components,
consistent with the long-range nature of the Coulomb interaction.

Both scalar perturbations~(\ref{eq:density-wave-kick}) and
(\ref{eq:Coulomb-pert}) share the spatial factor
$e^{\imath\mathbf{q}\cdot\mathbf{r}}$ with the EM
interaction~(\ref{eq:Hint-exact}), but the physical content is
different: the EM coupling probes the \emph{current--current}
response $\sigma(\omega)$ and requires gauge-consistent treatment,
while the scalar perturbations probe the \emph{density--density}
response $\chi(\mathbf{q},\omega)$ with no gauge involvement.
The extension of the present \texttt{RMG} implementation to
scalar-potential perturbations~(\ref{eq:density-wave-kick}) is
deferred to future work; in this paper we restrict to the
electromagnetic coupling in the long-wavelength limit throughout.


%--------------------------------------------------------------------
\subsection{Gauge Transformation: Length to Velocity Gauge}
\label{sec:gauge}
%--------------------------------------------------------------------

%....................................................................
\subsubsection{Why the length gauge fails for periodic systems}
\label{sec:LG-failure}
%....................................................................

In the long-wavelength limit the two standard gauge choices for the
electromagnetic coupling are:
\begin{alignat}{2}
  &\text{Length gauge (LG):}\quad
  &&\Avec^\mathrm{LG} = \mathbf{0},\quad
  \phi^\mathrm{LG}(\mathbf{r},t) = -\hat{\mathbf{r}}\cdot\Evec(t),
  \label{eq:LG}
  \\
  &\text{Velocity gauge (VG):}\quad
  &&\Avec^\mathrm{VG}(t)
    = -\!\int_{-\infty}^{t}\!\Evec(t')\,dt',\quad
  \phi^\mathrm{VG} = 0.
  \label{eq:VG}
\end{alignat}
In the length gauge the interaction Hamiltonian is
$\hat{HH_HL\mathrm{int}^\mathrm{LG} = \hat{\mathbf{r}}\cdot\Evec(t)$,
which requires evaluating matrix elements
$\langle\psi_a|\hat{rH_HL\alpha|\psi_b\rangle$ of the position
operator between the propagated Kohn-Sham orbitals.

The motivation for the gauge transformation from LG to VG is
precisely to eliminate $\hat{\mathbf{r}}$ from the interaction.
For a finite molecular system, as in paper~1,
\cite{jakowski-rmg-tddft-2025}
the position operator is well-defined and the length gauge is the
natural choice.
For a periodic system with Born--von K\'{a}rm\'{a}n boundary
conditions, the position operator is incompatible with the periodicity
of the Bloch states: the expectation value
$\langle\psi_{n\mathbf{k}}|\hat{rH_HL\alpha|\psi_{m\mathbf{k}}\rangle$
is not translationally invariant and is ill-defined for extended
states.\cite{resta-1998,king-smith-vanderbilt-1993}
The velocity gauge eliminates this problem by replacing $\hat{\mathbf{r}}$
with the momentum operator $\hat{\mathbf{p}}$, whose matrix elements
between Bloch states are well-defined.
We emphasize that this gauge motivation is specific to the
electromagnetic coupling: the scalar-potential perturbations of
Section~\ref{sec:scalar-pert} do not contain $\hat{\mathbf{r}}$
and are free of this difficulty for any $\mathbf{q}\neq\mathbf{0}$.


%....................................................................
\subsubsection{The G\"{o}ppert-Mayer gauge transformation}
\label{sec:GM}
%....................................................................

The length and velocity gauges are connected by the gauge function
\begin{equation}
  \chi(\mathbf{r},t) = -\hat{\mathbf{r}}\cdot\Avec^\mathrm{VG}(t),
  \label{eq:chi}
\end{equation}
which generates the G\"{o}ppert-Mayer
transformation.\cite{goeppert-mayer-1931}
Under a gauge transformation with function $\chi$, the potentials
transform according to
\begin{equation}
  \Avec\to\Avec+\nabla\chi,
  \qquad
  \phi\to\phi-\frac{\partial\chi}{\partial t},
  \label{eq:gauge-rules}
\end{equation}
and the wavefunction acquires a local phase,
\begin{equation}
  \Psi(\mathbf{r},t)
  \to \Psi'(\mathbf{r},t)
  = \hat{U}(\mathbf{r},t)\,\Psi(\mathbf{r},t),
  \qquad
  \hat{U} = e^{-\imath\chi(\mathbf{r},t)}
           = e^{\imath\hat{\mathbf{r}}\cdot\Avec(t)},
  \label{eq:psi-gt}
\end{equation}
where $\hat{U}$ is unitary ($|\Psi'|^2=|\Psi|^2$: charge density
is gauge invariant).
The Hamiltonian transforms as
\begin{equation}
  \hat{H}
  \to \hat{H}'
  = \hat{U}\hat{H}\hat{U}^\dagger
    - \imath\hat{U}\frac{\partial\hat{U}^\dagger}{\partial t}.
  \label{eq:H-gt}
\end{equation}


%....................................................................
\subsubsection{The velocity-gauge Kohn-Sham Hamiltonian}
\label{sec:H-VG}
%....................................................................

We apply the transformation~(\ref{eq:H-gt}) to the length-gauge
Hamiltonian
$H_{LG} = -\nabla^2/2 + V_\mathrm{KS} + \Vnl
+ \hat{\mathbf{r}}\cdot\Evec(t)$,
evaluating each term in turn.

\paragraph{Kinetic energy.}
Acting on a test function $g(\mathbf{r})$:
\begin{align}
  \hat{U}\hat{pH_HL\alpha\hat{U}^\dagger\,g
  &= e^{\imath\mathbf{r}\cdot\Avec}
     (-\imath\nabla_\alpha)
     \!\left(e^{-\imath\mathbf{r}\cdot\Avec}\,g\right)
  \nonumber\\
  &= e^{\imath\mathbf{r}\cdot\Avec}
     (-\imath)
     \!\left[(-\imath A_\alpha)
       e^{-\imath\mathbf{r}\cdot\Avec}\,g
     + e^{-\imath\mathbf{r}\cdot\Avec}\nabla_\alpha g
     \right]
   = (\hat{pH_HL\alpha - A_\alpha)\,g,
  \label{eq:U-p}
\end{align}
so $\hat{U}\hat{\mathbf{p}}\hat{U}^\dagger = \hat{\mathbf{p}}-\Avec$
and the kinetic energy transforms as
$\hat{U}(-\nabla^2/2)\hat{U}^\dagger
= \frac{1}{2}(\hat{\mathbf{p}}-\Avec)^2$.
For an electron with charge $-1$, carrying the sign through the
minimal substitution gives the VG kinetic term
$\frac{1}{2}(\hat{\mathbf{p}}+\Avec)^2$ (see footnote in
Section~\ref{sec:plane-wave}).

\paragraph{Time-derivative term.}
\begin{equation}
  -\imath\hat{U}\frac{\partial\hat{U}^\dagger}{\partial t}
  = -\imath\,e^{\imath\mathbf{r}\cdot\Avec}
    \frac{\partial}{\partial t}
    e^{-\imath\mathbf{r}\cdot\Avec}
  = -\mathbf{r}\cdot\dot{\Avec}(t)
  = +\mathbf{r}\cdot\Evec(t),
  \label{eq:dt-term}
\end{equation}
where $\Evec=-\partial\Avec/\partial t$.
This term $+\mathbf{r}\cdot\Evec$ exactly cancels the length-gauge
interaction $+\mathbf{r}\cdot\Evec$ in $\H_{LG}$:
the G\"{o}ppert-Mayer transformation is designed precisely so that
this cancellation occurs.

\paragraph{Local potential.}
$V_\mathrm{KS}(\mathbf{r})$ commutes with $\hat{U}(\mathbf{r},t)$
since both are diagonal in position space:
$\hat{U}V_\mathrm{KS}\hat{U}^\dagger = V_\mathrm{KS}$.

\paragraph{Nonlocal pseudopotential.}
The nonlocal pseudopotential $\Vnl$ is not diagonal in position
space, so $[\hat{U},\Vnl]\neq0$.
Expanding to first order in $A_0$:
\begin{equation}
  \hat{U}\Vnl\hat{U}^\dagger
  = e^{\imath\hat{\mathbf{r}}\cdot\Avec}\Vnl
    e^{-\imath\hat{\mathbf{r}}\cdot\Avec}
  \approx \Vnl
  + \imath\Avec(t)\cdot[\hat{\mathbf{r}},\Vnl]
  + \mathcal{O}(A_0^2).
  \label{eq:Vnl-gt}
\end{equation}
The correction $\imath\Avec\cdot[\hat{\mathbf{r}},\Vnl]$ is
identical to the Starace commutator that appears in the velocity
operator (Section~\ref{sec:current}): the gauge-transformation
and Heisenberg-equation-of-motion routes to the nonlocal
correction are equivalent, providing an internal consistency
check.\cite{starace-1971,Mattiat-Luber-2022}
The full algebra for the Kleinman-Bylander form of $\Vnl$ is
given in Appendix~\ref{app:gauge_nlpp}.

\paragraph{Full velocity-gauge Hamiltonian.}
Collecting all contributions and noting the cancellation of the
length-gauge coupling:
\begin{equation}
  \boxed{
  \HVG(t)
  = \frac{1}{2}\!\left(\hat{\mathbf{p}}+\Avec(t)\right)^{\!2}
  + V_\mathrm{KS}(\mathbf{r},t)
  + \Vnl
  + \imath\Avec(t)\cdot\!\left[\hat{\mathbf{r}},\Vnl\right],
  }
  \label{eq:H-VG}
\end{equation}
which contains no position operator $\hat{\mathbf{r}}$ in the
interaction and is therefore well-defined for periodic
systems.\cite{Mattiat-Luber-2022,pickard-mauri-2001}
For a periodic system with Bloch ansatz
$\psi_{n\mathbf{k}}=e^{\imath\mathbf{k}\cdot\mathbf{r}}u_{n\mathbf{k}}$,
the gradient acting on the Bloch factor contributes an additional
$+\mathbf{k}$ shift to the canonical momentum
(Section~\ref{sec:current-bloch}), giving the $\mathbf{k}$-dependent
form
\begin{equation}
  \hat{H}^\mathrm{VGH_HL\mathbf{k}(t)
  = \frac{1}{2}
    \!\left(-\imath\nabla+\mathbf{k}+\Avec(t)\right)^{\!2}
  + V_\mathrm{KS}(\mathbf{r},t)
  + \Vnl
  + \imath\Avec(t)\cdot\!\left[\hat{\mathbf{r}},\Vnl\right].
  \label{eq:H-VG-k}
\end{equation}


%....................................................................
\subsubsection{The perturbative momentum kick}
\label{sec:delta-kick}
%....................................................................

In the linear response regime the vector potential is applied as a
single instantaneous impulse at $t=0$:
\begin{equation}
  \Avec(t) = A_0\,\hat{\varepsilon}\,\delta(t),
  \label{eq:delta-kick}
\end{equation}
corresponding to a constant spectral amplitude
$\widetilde{\Evec}(\omega)=A_0\hat{\varepsilon}$ for all
frequencies.\cite{yabana-bertsch-1996}
After the kick ($t>0$), $\Avec=0$ and $\HVG$ reduces to
$\hat{HH_HL0+V_\mathrm{KS}[\rho(t)]$; the subsequent time evolution
is driven entirely by the self-consistent density response.
In the small-amplitude limit $|A_0|\ll1$, the kick imparts a
uniform phase to each orbital,
\begin{equation}
  \psi_{n\mathbf{k}}(\mathbf{r},0^+)
  \approx
  \left(1+\imath A_0\hat{\varepsilon}\cdot\hat{\mathbf{p}}\right)
  \psi_{n\mathbf{k}}^{(0)}(\mathbf{r}),
  \label{eq:kick-phase}
\end{equation}
exciting all frequencies simultaneously.
The full optical conductivity tensor
$\sigma_{\alpha\beta}(\omega)$ is recovered from three independent
simulations with kick directions
$\hat{\varepsilon}\in\{\hat{x},\hat{y},\hat{z}\}$,
as described in Section~\ref{sec:postproc}.


%--------------------------------------------------------------------
\subsection{Vector Potential and Magnetic Perturbations}
\label{sec:vector-potential}
%--------------------------------------------------------------------

The velocity-gauge Hamiltonian~(\ref{eq:H-VG}) provides a unified
framework for both electric and magnetic perturbations, distinguished
solely by the spatial structure of $\Avec(\mathbf{r},t)$.

%....................................................................
\subsubsection{Electric perturbation: spatially uniform \texorpdfstring{$\Avec(t)$}{A(t)}}
\label{sec:electric}
%....................................................................

In the long-wavelength approximation, $\Avec(\mathbf{r},t)\to\Avec(t)$
is spatially uniform.
The magnetic field of a spatially uniform vector potential vanishes
identically,
\begin{equation}
  \Bvec(t) = \nabla\times\Avec(t) = \mathbf{0},
  \label{eq:B-zero}
\end{equation}
so the perturbation is purely electric.
This is the regime of the present work.

%....................................................................
\subsubsection{Magnetic perturbation: spatially varying
  \texorpdfstring{$\Avec(\mathbf{r})$}{A(r)}}
\label{sec:magnetic}
%....................................................................

A static or time-dependent magnetic field $\Bvec$ is represented by
a spatially varying vector potential.
In the symmetric gauge centered at origin $\mathbf{RH_HLG$:
\begin{equation}
  \Avec(\mathbf{r})
  = \tfrac{1}{2}\Bvec\times(\mathbf{r}-\mathbf{RH_HLG),
  \label{eq:sym-gauge}
\end{equation}
from which $\nabla\times\Avec=\Bvec$ follows directly.
The orbital magnetic coupling
$\Avec(\mathbf{r},t)\cdot\hat{\mathbf{p}}$ in eq.~(\ref{eq:H-VG})
is then spatially dependent and drives magnetically induced transitions.
This framework is the theoretical foundation for computing
magnetically induced circular dichroism (ECD/MCD) spectra from
first principles.\cite{Mattiat-Luber-2022}
The gauge origin $\mathbf{RH_HLG$ in eq.~(\ref{eq:sym-gauge}) introduces
a formal ambiguity: for a complete basis physical observables are
independent of $\mathbf{RH_HLG$, while for a finite basis residual
dependence appears and vanishes in the limit of fine grid spacing.
\myblue{[cite: gauge-origin discussion, Mattiat Luber 2022 and
Pickard Mauri 2001 for real-space grids]}

%....................................................................
\subsubsection{Scope of the present implementation}
\label{sec:scope}
%....................................................................

All calculations in this work employ the electric perturbation of
Section~\ref{sec:electric}: a spatially uniform delta-kick vector
potential~(\ref{eq:delta-kick}) with $\Bvec=0$ for all $t>0$.
The diamagnetic term $\Avec^2/2$ in eq.~(\ref{eq:H-expand}) and the
diamagnetic current $\Jdia=\Avec|\psi|^2$ both vanish identically
after the kick (Section~\ref{sec:para-dia}), consistent with the
linear-response regime.
Extensions to the magnetic perturbation~(\ref{eq:sym-gauge}) for
ECD and MCD spectroscopy, and to the scalar-potential
perturbation~(\ref{eq:density-wave-kick}) for finite-wavevector
response functions, are natural next steps within the same
implementation framework and are left for future work.
